% =============================================================================
% Rozdział 7: Podsumowanie
% =============================================================================
\chapter{Podsumowanie}

\section{Podsumowanie zrealizowanych prac}

W ramach niniejszej pracy magisterskiej zrealizowano następujące cele:

\begin{enumerate}
  \item \textbf{Analiza literatury} — przeprowadzono przegląd istniejących badań nad autoskalowaniem w \textit{Kubernetes}, identyfikując pięć kluczowych luk badawczych: brak systematycznego porównania strategii HPA dla różnych profili obciążeniowych, ograniczoną analizę HPA Memory, brak badań łączących HPA z modelem \textit{SaaS}, brak analizy zjawisk anomalnych w autoskalowaniu chmurowym oraz brak otwartego, odtwarzalnego środowiska badawczego.

  \item \textbf{Zaprojektowanie architektury} — zaprojektowano system badawczy oparty na macierzy $2 \times 3$ (dwa profile obciążeniowe $\times$ trzy strategie HPA), składający się z czterech mikroserwisów (\texttt{cpu-service}, \texttt{io-service}, \texttt{echo-server}, \texttt{memory-service}), infrastruktury monitorującej opartej na \textit{kube-prometheus-stack} (Helm) oraz dedykowanej maszyny \textit{Compute Engine} jako generatora obciążenia.

  \item \textbf{Implementacja} — zaimplementowano kompletny system badawczy w języku Python 3.12 z frameworkiem \textit{FastAPI}, konteneryzowany za pomocą \textit{Docker} i wdrażany na zarządzanym klastrze \textit{Google Kubernetes Engine} (GKE). Wszystkie manifesty \textit{Kubernetes} (Deploymenty, Serwisy typu \texttt{LoadBalancer}, HPA) oraz konfiguracje monitoringu są dostępne w repozytorium. Serwisy są eksponowane przez \textit{GCP LoadBalancer}, co odzwierciedla warunki produkcyjne.

  \item \textbf{Automatyzacja testów} — opracowano skrypt \texttt{run-tests.sh} orkiestrujący wykonanie 8 scenariuszy testowych z 5-krotnym powtórzeniem każdego, z automatycznym przełączaniem HPA, resetowaniem \textit{Deploymentów} i cooldownem między scenariuszami. Skrypt działa na maszynie \texttt{k6-generator} (\textit{Compute Engine}) w tej samej sieci VPC co klaster GKE.

  \item \textbf{Scenariusze obciążeniowe} — zdefiniowano trzy skrypty \textit{k6} (\texttt{cpu-load-test.js}, \texttt{io-load-test.js}, \texttt{combined-scenario.js}) z precyzyjnymi progami akceptacji, profilami obciążenia dopasowanymi do charakterystyki aplikacji oraz mapowaniem hostów (\texttt{options.hosts}) na adresy IP \textit{GCP LoadBalancer}, eliminującym potrzebę konfiguracji DNS.

  \item \textbf{Infrastruktura monitorująca} — wdrożono \textit{kube-prometheus-stack} przez \textit{Helm}, zapewniając automatycznie skonfigurowany stos monitorujący (\textit{Prometheus}, \textit{Grafana}, \textit{kube-state-metrics}, \textit{node-exporter}) z \textit{Grafaną} eksponowaną publicznie przez \texttt{LoadBalancer}. Zdefiniowano procedurę eksportu 5 zestawów danych CSV na scenariusz (\texttt{replicas}, \texttt{cpu\_mcpu}, \texttt{memory\_mb}, \texttt{rps}, \texttt{latency\_ms}), które wraz z wynikami \textit{k6} tworzą kompletny zbiór danych badawczych.

  \item \textbf{Charakterystyka zjawisk anomalnych} — zidentyfikowano i opisano trzy kluczowe zjawiska przejściowe wpływające na autoskalowanie w środowisku chmurowym: \textit{HPA Lag} (opóźnienie reakcji autoskalera, $\sim$60--90\,s), \textit{Cold Start Latency} (opóźnienie zimnego startu Podów, $\sim$15--30\,s) oraz \textit{CPU Throttling Cascade} (kaskadowe błędy przy throttlingu CPU). Dla każdego zjawiska zaproponowano metodę ilościowej analizy na podstawie danych z eksportów CSV z \textit{Grafany}.
\end{enumerate}

\section{Weryfikacja hipotezy badawczej}

% TODO: uzupełnić po zakończeniu badań eksperymentalnych

Hipoteza badawcza postawiona w Rozdziale 1 brzmiała:

\begin{quote}
\textit{,,Nie istnieje jedna uniwersalna metryka HPA — skuteczność strategii skalowania zależy od charakterystyki obciążeniowej aplikacji (\textit{workload profile}). Błędny wybór metryki prowadzi do niedoskalowania lub przeskalowania.''}
\end{quote}

Na obecnym etapie prac (czerwiec 2026) badania eksperymentalne na klastrze GKE nie zostały jeszcze przeprowadzone, w związku z czym hipoteza nie może zostać formalnie zweryfikowana. Poniżej przedstawiono oczekiwane rezultaty oparte na analizie teoretycznej i charakterystyce zaimplementowanych aplikacji:

\begin{itemize}
  \item \textbf{HPA CPU} — oczekuje się, że będzie skuteczne wyłącznie dla \texttt{cpu-service}, gdzie zużycie procesora silnie koreluje z liczbą żądań (operacje \textit{Pillow}, rekurencyjny Fibonacci). Dla \texttt{io-service}, gdzie procesor jest bezczynny podczas \texttt{asyncio.sleep()}, HPA CPU nie powinno reagować na wzrost ruchu.

  \item \textbf{HPA Memory} — oczekuje się, że będzie nieskuteczne dla obu serwisów, ponieważ ani \texttt{cpu-service}, ani \texttt{io-service} nie alokują pamięci dynamicznie w odpowiedzi na obciążenie. Pamięć jest alokowana statycznie przy starcie i utrzymuje stały poziom.

  \item \textbf{HPA Custom RPS} — oczekuje się, że będzie uniwersalnie skuteczne, ponieważ liczba żądań na sekundę rośnie proporcjonalnie do obciążenia niezależnie od tego, czy żądania są CPU-intensywne, czy I/O-intensywne.
\end{itemize}

Dodatkowo, oczekuje się, że analiza zjawisk anomalnych (\textit{HPA Lag}, \textit{Cold Start Latency}, \textit{CPU Throttling Cascade}) dostarczy ilościowych danych na temat wpływu tych zjawisk na wydajność systemów \textit{SaaS} w środowisku GKE.

\section{Ograniczenia pracy}

Należy wskazać następujące ograniczenia niniejszej pracy:

\begin{itemize}
  \item \textbf{Brak wyników eksperymentalnych} — na obecnym etapie zaimplementowano kompletny system badawczy i narzędzia do testów, ale same badania nie zostały jeszcze przeprowadzone. Wnioski przedstawione w Rozdziale 5 mają charakter oczekiwany i wymagają walidacji empirycznej.

  \item \textbf{Sztuczny charakter obciążenia} — aplikacje badawcze (\texttt{cpu-service}, \texttt{io-service}) symulują określone profile obciążeniowe, ale nie odzwierciedlają rzeczywistych aplikacji produkcyjnych. Rzeczywiste systemy często wykazują mieszane profile (\textit{CPU-bound} + \textit{I/O-bound}).

  \item \textbf{Ograniczony zakres HPA} — badanie ogranicza się do trzech metryk HPA (\textit{CPU}, \textit{Memory}, \textit{Custom RPS}). Nie analizowano innych strategii, takich jak HPA na podstawie liczby żądań w toku (\textit{in-flight}), metryk zewnętrznych (\textit{External Metrics}) ani skalowania predykcyjnego.

  \item \textbf{Pojedyncza strefa GCP} — wszystkie eksperymenty są przeprowadzane w jednej strefie GCP (\texttt{europe-central2}), co oznacza, że nie badano wpływu opóźnień między strefami (\textit{cross-zone latency}) ani scenariuszy \textit{multi-region}.

  \item \textbf{Brak implementacji \textit{multi-tenancy}} — rozszerzenie \textit{SaaS Multi-Tenancy} (Faza VIII) zostało zaprojektowane koncepcyjnie, ale nie zaimplementowane w kodzie. Wnioski dotyczące zachowania systemu w środowisku wielodostępnym pozostają hipotetyczne.

  \item \textbf{Stałe progi HPA} — progi HPA (CPU 50\%, Memory 70\%, RPS 100/Pod) zostały ustalone arbitralnie na podstawie wartości referencyjnych. Nie badano wpływu różnych wartości progowych na skuteczność skalowania.
\end{itemize}

\section{Rekomendacje praktyczne}

Na podstawie analizy teoretycznej i projektu systemu sformułowano następujące rekomendacje dla praktyków:

\begin{enumerate}
  \item \textbf{Dobieraj metrykę HPA do profilu aplikacji} — przed wyborem strategii autoskalowania należy scharakteryzować aplikację pod kątem profilu obciążeniowego (\textit{CPU-bound}, \textit{I/O-bound}, \textit{memory-bound}). Domyślna strategia HPA CPU może być nieskuteczna dla aplikacji I/O-bound.

  \item \textbf{Używaj HPA Custom RPS jako strategii uniwersalnej} — jeśli charakterystyka aplikacji jest nieznana lub zmienna, metryka RPS (dostarczana przez \textit{prometheus-adapter}) jest bezpiecznym wyborem, ponieważ koreluje z obciążeniem niezależnie od profilu.

  \item \textbf{Unikaj HPA Memory jako jedynej metryki} — chyba że aplikacja wykazuje silną korelację między obciążeniem a zużyciem pamięci (np. cache, kolejki), HPA Memory prawdopodobnie nie będzie skuteczne.

  \item \textbf{Uwzględnij zjawiska anomalne w SLA} — \textit{HPA Lag} ($\sim$60--90\,s) i \textit{Cold Start Latency} ($\sim$15--30\,s) oznaczają, że system potrzebuje 1--2 minut na reakcję na nagły wzrost obciążenia. Progi SLA powinny być definiowane z uwzględnieniem tego opóźnienia — oczekiwanie, że system zareaguje w kilka sekund na gwałtowny skok ruchu, jest nierealistyczne przy HPA opartym na metrykach z \textit{Prometheusa}.

  \item \textbf{Monitoruj CPU throttling} — \textit{CPU throttling} jest ,,cichym zabójcą'' wydajności — nie powoduje awarii, ale drastycznie wydłuża czasy odpowiedzi. Rekomenduje się monitorowanie metryki \texttt{container\_cpu\_cfs\_throttled\_seconds\_total} i ustawienie alertów przy jej wzroście.

  \item \textbf{Używaj GCP LoadBalancer zamiast port-forward} — dla testów obciążeniowych i środowisk produkcyjnych, bezpośrednia ekspozycja serwisów przez \textit{GCP LoadBalancer} eliminuje wąskie gardło tunelowania przez \textit{API server} i zapewnia realistyczne warunki sieciowe.

  \item \textbf{Wdróż monitoring przez Helm} — \textit{kube-prometheus-stack} przez \textit{Helm} redukuje złożoność konfiguracji monitoringu i zapewnia spójność wersji komponentów, co jest szczególnie istotne w środowiskach produkcyjnych.
\end{enumerate}

\section{Dalsze kierunki badań}

Niniejsza praca otwiera następujące kierunki dalszych badań:

\begin{enumerate}
  \item \textbf{Przeprowadzenie badań eksperymentalnych na GKE} — najpilniejszym kolejnym krokiem jest wykonanie 8 scenariuszy testowych na klastrze GKE i zebranie danych empirycznych do weryfikacji hipotezy. Środowisko badawcze, narzędzia i procedury są w pełni przygotowane.

  \item \textbf{Ilościowa analiza zjawisk anomalnych} — wykorzystanie 5 eksportów CSV na scenariusz do ilościowego scharakteryzowania \textit{HPA Lag}, \textit{Cold Start Latency} i \textit{CPU Throttling Cascade}. Dane te mogą posłużyć do opracowania modeli predykcyjnych wykrywających te zjawiska.

  \item \textbf{Implementacja \textit{multi-tenancy}} — zrealizowanie Fazy VIII zgodnie z koncepcją opisaną w Rozdziale 6 i przeprowadzenie testów w środowisku wielodostępnym na GKE.

  \item \textbf{Rozszerzenie macierzy badawczej} — dodanie \texttt{memory-service} jako trzeciego profilu obciążeniowego (\textit{memory-bound}) i zbadanie, czy HPA Memory staje się skuteczne dla aplikacji, których zużycie pamięci koreluje z obciążeniem.

  \item \textbf{Badanie innych metryk HPA} — analiza skuteczności metryk takich jak liczba żądań w toku (\texttt{http\_requests\_inflight}), metryk zewnętrznych (długość kolejki \textit{RabbitMQ}, liczba wiadomości \textit{Kafka}) oraz kombinacji wielu metryk w jednym HPA.

  \item \textbf{Skalowanie predykcyjne} — porównanie reaktywnego HPA z podejściami predykcyjnymi (np. opartymi na LSTM), które przewidują przyszłe obciążenie i skalują proaktywnie, potencjalnie eliminując problem \textit{HPA Lag}.

  \item \textbf{Analiza kosztowa} — rozszerzenie badań o analizę kosztów infrastruktury GCP w modelu \textit{pay-as-you-go} dla różnych strategii HPA, szczególnie w kontekście \textit{multi-tenancy}, gdzie przeskalowanie przekłada się bezpośrednio na wyższe koszty.

  \item \textbf{Testy wielostrefowe} — powtórzenie eksperymentów w konfiguracji \textit{multi-zone} GKE w celu zbadania wpływu opóźnień między strefami na skuteczność HPA i zjawiska anomalne.

  \item \textbf{Testy z rzeczywistymi aplikacjami} — zastąpienie syntetycznych serwisów badawczych rzeczywistymi aplikacjami (np. system rezerwacji, przetwarzanie obrazu w pipeline ML) w celu walidacji uniwersalności wniosków.
\end{enumerate}

\section{Podsumowanie końcowe}

Niniejsza praca magisterska stanowi kompleksowe przygotowanie do badań empirycznych nad skutecznością strategii \textit{HorizontalPodAutoscaler} w \textit{Kubernetes} dla różnych profili obciążeniowych, w całości w środowisku \textit{Google Cloud Platform}. Zaimplementowano w pełni funkcjonalny system badawczy składający się z mikroserwisów, infrastruktury monitorującej opartej na \textit{kube-prometheus-stack} (Helm) oraz zestawu skryptów automatyzujących, gotowy do przeprowadzenia eksperymentów na klastrze GKE.

Główny wkład pracy obejmuje:

\begin{itemize}
  \item Systematyczną analizę literatury identyfikującą luki badawcze w dziedzinie autoskalowania w \textit{Kubernetes}, w tym nowo zidentyfikowaną lukę dotyczącą analizy zjawisk anomalnych (\textit{HPA Lag}, \textit{Cold Start Latency}, \textit{CPU Throttling Cascade}).
  \item Projekt macierzy badawczej $2 \times 3$ umożliwiającej izolowane porównanie strategii HPA.
  \item Implementację czterech mikroserwisów o kontrolowanych profilach obciążeniowych (\textit{CPU-bound}, \textit{I/O-bound}, \textit{memory-bound}).
  \item Wdrożenie infrastruktury monitorującej przez \textit{kube-prometheus-stack} (Helm) z niestandardowymi metrykami przez \textit{prometheus-adapter}.
  \item Pełną automatyzację procesu badawczego — od wdrożenia po wykonanie testów i eksport danych CSV z \textit{Grafany}.
  \item Koncepcję rozszerzenia \textit{SaaS Multi-Tenancy} dla środowiska GKE, uwzględniającą specyfikę zarządzanego klastra i mechanizmy izolacji \textit{Kubernetes}.
  \item Otwarte, odtwarzalne środowisko badawcze dostępne w repozytorium, możliwe do odtworzenia na dowolnym klastrze GKE.
\end{itemize}

Przeprowadzenie zaplanowanych badań eksperymentalnych i weryfikacja hipotezy badawczej stanowią bezpośredni kolejny krok, do którego niniejsza praca dostarcza wszystkich niezbędnych narzędzi. Szczególnie wartościowym wkładem będzie ilościowa analiza zjawisk anomalnych na podstawie danych z 5 eksportów CSV na scenariusz — wyniki te mogą mieć istotne znaczenie praktyczne dla projektantów systemów \textit{SaaS} działających na \textit{Kubernetes} w chmurze.
